\begin{table}[htbp]
\centering
\small
\caption[Karakteristik Zona Morfometri Transekt]{Karakteristik Gradien Presipitasi Sepanjang Profil Morfometri Transekt Utara--Selatan Kebumen}
\label{tab:orographic_transect_zones}
\resizebox{\linewidth}{!}{%
\begin{tabular}{lccccc}
\hline
\textbf{Zona Morfologi Transekt} & \textbf{Elevasi Rerata (m)} & \textbf{GSMaP (mm)} & \textbf{CHIRPS (mm)} & \textbf{Fused (mm)} & \textbf{Selisih (mm)} \\
\hline
Zona Pegunungan Utara (Sadang/Karangsambung) & 223 & 2597.2 & 3353.6 & 3018.9 & -756.4 \\
Zona Perbukitan Transisi (Alian/Pejagoan) & 209 & 2588.0 & 3336.5 & 3002.6 & -748.5 \\
Zona Dataran Aluvial Tengah (Kebumen/Kutowinangun) & 48 & 2533.9 & 3256.0 & 2930.0 & -722.1 \\
Zona Pesisir Pantai Selatan (Samudera Hindia) & 13 & 2539.8 & 3265.5 & 2931.8 & -725.7 \\
\hline
\end{tabular}%
}
\end{table}
